\section{Условие}
{\bfseries Цель работы:}
Приобретение практических навыков в управлении потоками в операционной системе и обеспечении синхронизации между потоками.

{\bfseries Задание:}
Составить программу на языке Си, обрабатывающую данные в многопоточном режиме. При обработке использовать стандартные средства создания потоков операционной системы (Windows/Unix). Ограничение максимального количества потоков, работающих в один момент времени, должно быть задано ключом запуска вашей программы.

Также необходимо уметь продемонстрировать количество потоков, используемое вашей программой с помощью стандартных средств операционной системы.

В отчёте привести исследование зависимости ускорения и эффективности алгоритма от входных данных и количества потоков. Получившиеся результаты необходимо объяснить.

{\bfseries Вариант:} 16

{\bfseries Описание варианта 16:}
Задаётся радиус окружности. Необходимо с помощью метода Монте-Карло рассчитать её площадь.

{\bfseries Метод Монте-Карло для вычисления площади окружности:}
\begin{itemize}
    \item Окружность радиуса R вписывается в квадрат со стороной 2R
    \item Генерируется N случайных точек, равномерно распределённых в квадрате
    \item Подсчитывается количество точек M, попавших внутрь окружности
    \item Площадь окружности вычисляется по формуле: $S = \frac{M}{N} \times S_{квадрата} = \frac{M}{N} \times 4R^2$
    \item Оценка числа $\pi$: $\pi \approx \frac{M}{N} \times 4$
\end{itemize}
